\documentclass{plantillaPPT}

\titulo{8}{Volúmenes de Sólidos de Revolución}

\begin{document}

\primeraDiapo

\begin{frame}{Método de Discos o Anillos}
    \framesubtitle{Métodos para obtener volúmenes de sólidos de revolución}

Expresión general:

\[V = \int_{a}^{b} \pi \, r^2\; dx \qquad\text{o bien}\qquad V=\int_{a}^{b}\pi (R^2-r^2)dx\]

Con \(R\) y \(r\) los radios del disco/anillo a generar.\\ \vspace{0.2cm}
\begin{itemize}
    \item Si se usa la expresión 2, \(R\) es el radio mayor y \(r\) es el radio menor del anillo. \(\Rightarrow R>r\)
    \item Se integra de forma \textbf{paralela} al eje de rotación.
    \item Los radios son segmentos \textbf{perpendiculares} al eje de rotación.

\end{itemize}

\end{frame}

\begin{frame}{Método de Capas Cilíndricas}
    \framesubtitle{Métodos para obtener volúmenes de sólidos de revolución}

Expresión general:
\[\int_{a}^{b} 2\pi \,r\, h \; dx\]
Con \textbf{\(r\)} el \textbf{radio} y \textbf{\(h\)} la \textbf{altura} de las capas a generar. \\ \vspace{0.2cm}
\begin{itemize}
    \item Se integra de forma \textbf{perpendicular} al eje de rotación.
    \item Los radios son segmentos \textbf{perpendiculares} al eje de rotación. 
    \item Las alturas son segmentos \textbf{paralelos} al eje de rotación.
\end{itemize}
\end{frame}

\begin{frame}{Volúmenes de Sólidos de Revolución}
    \framesubtitle{Práctica 8}

2. Sea \(R\) la región definida por:
\[R = \{(x,y)\in \mathbb{R}^2:0\leq x\leq 2\pi \land 0 \leq y \leq |\sin(x)|\}\]
Escriba la o las integrales que permiten calcular el volumen del sólido que se genera al hacer rotar R en torno a los ejes:
\[(a) \; y = 0 \qquad \qquad (b)\;x=0 \]
    
\end{frame}

\begin{frame}{Condiciones para búsqueda}
    \framesubtitle{Práctica 8}
    
2. Sea \(R\) la región del plano limitada por la parábola \(y=x^2\) y la recta \(y=mx\), con \(m>0\). Determinar el valor de \(m\) de modo que los volúmenes generados por la rotación de \(R\) en torno al eje \(X\) y en torno al eje \(Y\) sean iguales.
    
\end{frame}

\begin{frame}{Volúmenes de Sólidos de Revolución}
    \framesubtitle{Práctica 8}

3. Sea \(R\) la región del plano acotada por las gráficas \(C_1:x=\ln(y)\) y \(C_2:y=\ln(x)\) entre las rectas \(y=1\) e \(y=2\). \\ \vspace{0.3cm}

\((a)\) Calcular el valor del área de la región \(R\). \\ \vspace{0.3cm}
\((b)\) Escriba la o las integrales que permiten calcular el volumen del sólido que se genera al hacer rotar \(R\) en torno a los ejes \(x=8\) \(y=4\). \\ \vspace{0.3cm}

\end{frame}

\begin{frame}{Secciones Transversales}
    \framesubtitle{Práctica 8}

\((c)\) Escriba la o las integrales que permiten calcular el volumen del sólido tal que su base es \(R\) y cuyas secciones transversales perpendiculares al eje \(Y\) son rectángulos para los cuales la altura es 3 veces la base.

\end{frame}

\end{document}